\chapter{Quand les marchés se brisent ensemble}\label{ch:m3:when-markets-break-together}

En plein mois de mars 2020, alors que la pandémie fermait les économies du monde, les actions chutèrent, et il en alla de même des actifs censés monter quand les actions baissent.
Entre le 9 et le 18 mars, le rendement des bons du Trésor américain à dix ans, l'actif le plus sûr du monde,
passa de 0.54~\% à 1.18~\% à mesure que les détenteurs les vendaient~; l'indice du dollar élargi
progressa de plus de 7~\%~; le 16 mars, le S\&P 500 perdit 12~\%, sa pire séance depuis 1987. Les investisseurs ne choisissaient plus
entre actifs risqués et actifs sûrs~; ils vendaient tout ce qu'ils pouvaient pour se procurer des liquidités.
Les chapitres précédents décrivaient les crises propres à chaque marché~: les prix de l'électricité en 2022, le nickel au London Metal
Exchange, les \omterm{def:m3:margin-liquidation-and-loss-allocation:cascade}{cascades de liquidations} sur crypto-actifs, FTX. Ce chapitre porte sur les journées où les crises ne font qu'une, et sur les liens entre marchés qui, invisibles en temps calme,
détenteurs communs, financements communs, règles de marge communes, deviennent la seule chose qui compte.
Il en expose les mécanismes, passe en revue
2008, mars 2020 et 2022 du point de vue des desks de ce livre, et construit un
simulateur de la spirale.

\section{Les mécanismes des tensions conjointes~: ventes forcées, spirales des pertes et des marges}

\begin{definition}[Vente forcée]\label{def:m3:when-markets-break-together:fire}
Une \emph{vente forcée}\index{vente forcée} est une vente imposée à un détenteur par un besoin de liquidités ou un appel de marge, à un prix
inférieur à celui que vaudrait l'actif pour des acheteurs non contraints, parce que les acheteurs naturels de cet actif vendent eux aussi.
\end{definition}

\begin{definition}[Spirale des pertes, spirale des marges]\label{def:m3:when-markets-break-together:spirals}
Une \emph{spirale des pertes}\index{spirale des pertes} est la boucle de rétroaction où les pertes d'un détenteur à effet de levier le forcent à vendre, ses ventes font baisser les prix, et les
prix plus bas lui infligent de nouvelles pertes, à lui comme aux autres détenteurs des mêmes actifs. Une \emph{spirale des marges}\index{spirale des marges} est la boucle de rétroaction
où la hausse de la volatilité relève les exigences de marge, les marges plus élevées imposent un désendettement, et les \omterm{def:m3:when-markets-break-together:fire}{ventes forcées} accroissent encore la volatilité.
\end{definition}

Brunnermeier et Pedersen ont donné leur nom aux deux spirales dans un modèle de la façon dont la liquidité de financement et la liquidité de marché s'entretiennent mutuellement, la
spirale de liquidité du livre~1, chapitre~31. Deux propriétés les transforment en événement inter-marchés. Les détenteurs à effet de levier possèdent de nombreux actifs, si bien qu'une perte sur l'un oblige
à vendre l'ensemble de leurs positions~; les marges comme les décotes sont fixées sur la volatilité récente de chaque actif, si bien qu'elles montent partout à la fois. L'actif qui est vendu
n'est pas celui qui a chuté, mais celui qui peut être vendu.

\begin{omfigure}
\centering
\begin{tikzpicture}[font=\small,
  s/.style={draw,thick,rounded corners=2pt,align=center,minimum height=0.9cm,text width=2.6cm},
  arr/.style={-{Stealth[length=2.2mm]},thick}]
\node[s,draw=omThm,fill=omThm!8] (sh) at (0,0) {choc~: un\\actif chute};
\node[s,draw=omDef,fill=omDef!8] (eq) at (3.8,1.1) {fonds propres des\\détenteurs baissent};
\node[s,draw=omExo,fill=omExo!8] (mg) at (3.8,-1.1) {volatilité et\\marges augmentent};
\node[s,draw=omProp,fill=omProp!8] (fs) at (7.6,0) {ventes forcées de\\\emph{tous} les actifs détenus};
\node[s,draw=gray,fill=gray!8] (px) at (11.2,0) {les prix baissent\\partout};
\draw[arr] (sh) -- (eq); \draw[arr] (sh) -- (mg); \draw[arr] (eq) -- (fs); \draw[arr] (mg) -- (fs); \draw[arr] (fs) -- (px);
\draw[arr,dashed] (px.north) to[bend right=30] node[above,font=\footnotesize] {spirale des pertes} (eq.north);
\draw[arr,dashed] (px.south) to[bend left=30] node[below,font=\footnotesize] {spirale des marges} (mg.south);
\end{tikzpicture}

\omcaption{Les deux spirales. Un choc sur un actif réduit les fonds propres des détenteurs à effet de levier et fait monter la volatilité mesurée et les marges~; les deux imposent des ventes de
tout ce que les détenteurs possèdent~; les ventes font baisser les prix et accroissent la volatilité partout, ce qui alimente les deux spirales en retour. Schéma indicatif.}\label{fig:m3:when-markets-break-together:spirals}
\end{omfigure}

\begin{proposition}[Le multiplicateur de vente forcée sans impact]\label{prop:m3:when-markets-break-together:multiplier}
Un détenteur dont les fonds propres sont inférieurs de $S$ à son exigence de marge, sur des positions au taux de marge $m$, doit vendre pour $S/m$ de ces positions afin de rétablir sa marge,
si les prix ne bougent pas~: chaque dollar vendu libère $m$ d'exigence et laisse les fonds propres inchangés. Avec impact sur les cours, chaque vente fait baisser les cours des
positions conservées et, avec elles, les fonds propres, et les \omterm{def:m3:when-markets-break-together:fire}{ventes forcées} par dollar de déficit initial dépassent $1/m$.
\end{proposition}

\begin{proof}
Vendre un dollar de position au prix de marché ne modifie pas les fonds propres et réduit l'exigence de $m$~; rétablir $S$ demande $S/m$. Avec impact, les positions vendues
font évoluer les prix au détriment des positions conservées, ce qui crée un nouveau déficit qu'il faut combler par de nouvelles ventes.
\end{proof}

\begin{definition}[Fuite vers la qualité, rupture des corrélations]\label{def:m3:when-markets-break-together:flight}
Une \emph{fuite vers la qualité}\index{fuite vers la qualité} est un déplacement des investisseurs hors des actifs risqués vers ceux perçus comme les plus sûrs, ce qui fait généralement monter les
prix des obligations d'État tandis que les autres actifs reculent. Une \emph{rupture des corrélations}\index{rupture des corrélations} est un changement brutal des corrélations
entre actifs en période de tension, généralement vers un, les \omterm{def:m3:when-markets-break-together:fire}{ventes forcées} les faisant évoluer ensemble, ce qui invalide les couvertures et les estimations de risque calibrées en temps calme.
\end{definition}

\section{2008}

La crise de 2008 sert de modèle~: des établissements à effet de levier se finançaient à court terme contre des titres dont les valeurs et les décotes évoluaient ensemble, et la
défaillance de l'un d'eux, Lehman Brothers, le 15 septembre 2008, a transformé les doutes sur chacun en une ruée sur tous. Les chambres de compensation évoquées dans le livre~1 ont tenu~:
Le service de swaps de LCH a géré le portefeuille de Lehman, soit 66\,390 transactions pour un notionnel de 9~billions de dollars, dans les limites de la marge qu'elle détenait, en le couvrant et en le mettant aux enchères
entre le 24 septembre et le 3 octobre, sans recourir à son fonds de défaillance. Les pertes sont venues d'ailleurs, dans les opérations non compensées et les opérations financées~: le financement à court terme cessait d'être renouvelé
et des décotes en hausse sur le collatéral de toute nature. Les
enseignements qui ont façonné le système d'après 2008, la compensation centrale, la marge pour les dérivés non compensés et les règles de liquidité bancaire, expliquent pourquoi la crise suivante a semblé
différente.

\section{Mars 2020}

\begin{dated}{2026-09}{La tourmente de mars 2020}\label{dat:m3:when-markets-break-together:2020}
Le rapport du Financial Stability Board distingue une fuite vers la sécurité, de la fin février au début de mars 2020, au cours de laquelle les investisseurs ont vendu des actifs risqués et
acheté des actifs sûrs, et une «~ruée vers la liquidité~» (dash for cash) à la mi-mars, au cours de laquelle ils ont vendu indistinctement actifs risqués et actifs relativement sûrs pour obtenir des liquidités~; les actions et les obligations d'État
reculèrent alors ensemble, la volatilité augmenta même pour les obligations d'État et pour l'or, et les corrélations se rompirent. Le 16 mars, le S\&P 500 perdit 12~\%, sa pire
séance depuis 1987, et le VIX atteignit son plus haut historique, à 83. Le dysfonctionnement du marché des Treasuries a été aggravé par les ventes de certains investisseurs non bancaires
à effet de levier et des détenteurs étrangers. Le Comité de Bâle, le CPMI et l'OICV ont signalé que les exigences de marge initiale des contreparties centrales avaient augmenté d'environ
300~milliards de dollars sur le seul mois de mars 2020, avec un pic d'appel de marge de variation de 140~milliards de dollars le 9 mars.
\end{dated}

\begin{omfigure}
\begin{tikzpicture}
\begin{axis}[width=11cm,height=5cm,axis y line*=left,
  xlabel={jour de négociation depuis le 3 février 2020},ylabel={rendement à 10 ans~\%},
  tick label style={font=\small},label style={font=\small},
  xmin=0,xmax=61,ymin=0.4,ymax=1.7,grid=major,grid style={gray!25},
  table/col sep=comma]
\fill[omExo!15] (axis cs:24,0.4) rectangle (axis cs:34,1.7);
\addplot[omThm,very thick] table[x=day,y=ust10y]{figdata/markets-3/28-when-markets-break-together/march2020.csv};\label{plot:m3:ust10}
 \node[font=\footnotesize,anchor=north,text=omExo] at (axis cs:29,1.68) {ruée vers la liquidité};
\end{axis}
\begin{axis}[width=11cm,height=5cm,axis y line*=right,axis x line=none,
  ylabel={indice du dollar élargi},ylabel near ticks,
  tick label style={font=\small},label style={font=\small},
  xmin=0,xmax=61,ymin=112,ymax=128,
  legend columns=2,legend style={font=\footnotesize,at={(0.5,-0.3)},anchor=north,draw=none,fill=none,/tikz/every even column/.append style={column sep=8pt}},
  table/col sep=comma]
\addlegendimage{/pgfplots/refstyle=plot:m3:ust10}\addlegendentry{rendement à 10 ans (gauche)}
\addplot[omDef,thick,dashed] table[x=day,y=usd]{figdata/markets-3/28-when-markets-break-together/march2020.csv};\addlegendentry{indice du dollar élargi (droite)}
\end{axis}
\end{tikzpicture}

\omcaption{Actifs sûrs en mars 2020~: le rendement du Trésor à dix ans est tombé à 0.54~\% le 9 mars, dans la fuite vers la sécurité, puis est remonté à 1.18~\% le 18 mars, lorsque
les bons du Trésor ont été vendus pour obtenir du cash (zone grisée, du 9 au 23 mars), tandis que l'indice du dollar élargi a progressé de 7.6~\% du 9 au 23 mars. Données~: séries FRED DGS10 et DTWEXBGS (Board of Governors of the Federal Reserve System).
}\label{fig:m3:when-markets-break-together:march2020}
\end{omfigure}

Les desks de ce livre ont observé la même semaine à travers des fenêtres différentes. Le desk de taux a vu les bons du Trésor vendus avec les actions et la base des contrats à terme se creuser fortement
(livre~2). Le desk de matières premières a vu le pétrole s'effondrer pour des raisons qui lui sont propres (les chapitres sur le stockage de ce livre, et le règlement-livraison négatif du WTI
le mois suivant) tandis que les marges sur tous les actifs augmentaient. Un desk crypto a vu le bitcoin vendu avec tout le reste, quelle que soit sa réputation de couverture~: son cours sur Coinbase,
tel que FRED l'enregistre, est passé de 7\,957 dollars le 11 mars à 4\,980 le 12 mars, soit 37~\% en une journée. Chaque desk a vu ses marges augmenter la même semaine, ce qui est la \omterm{def:m3:when-markets-break-together:spirals}{spirale des marges} vue de l'intérieur.

\section{2022~: marges énergétiques, nickel, gilts et crypto}

En 2022, il n'y a pas eu de krach unique, mais une succession de spirales sur différents marchés, chacune entraînée par les mêmes mécanismes en un lieu différent. Les prix européens du gaz et de l'électricité
ont tant augmenté (\cref{ch:m3:natural-gas-and-lng}) que la marge de trésorerie sur les couvertures par contrats à terme est devenue, pour les entreprises qui couvrent des positions physiques, le risque de liquidité
décrit au \cref{ch:m3:getting-access-commodities-and-power}. La BCE a constaté que la marge initiale sur les portefeuilles de matières premières était environ deux fois plus élevée à la mi-2022
qu'à la fin 2021, les matières premières représentant environ 2~\% du notionnel brut des produits dérivés de la zone euro mais plus de 20~\% de la marge initiale versée, et quatre banques
compensaient environ 85~\% des positions énergétiques négociées en bourse~; l'importateur de gaz Uniper avait tiré la totalité d'une ligne de crédit de 9~milliards d'euros auprès de la banque publique allemande
KfW au 29 août, et demandait 4~milliards de plus, car les restrictions d'approvisionnement en gaz et les prix faisaient grimper ses besoins de marge.
En mars, le nickel à la London Metal Exchange est passé d'environ 30\,000 à plus de 100\,000 dollars la tonne en une journée et une matinée, les appels de marge de la journée ont atteint
environ 7~milliards de dollars, et trois membres compensateurs n'ont pas réglé à temps, si bien que la bourse a suspendu la négociation et annulé les transactions d'une matinée
(\cref{ch:m3:metals}). En septembre, les stratégies de gestion adossée au passif des fonds de pension britanniques ont fait face aux appels de marge sur leurs positions en gilts à mesure que les rendements longs montaient
plus vite que jamais, ont vendu des gilts pour les payer, et ont poussé les rendements encore plus haut, jusqu'à ce que la Banque d'Angleterre achète des gilts longs (livre~2, chapitre~7). Dans la crypto, le
\omterm{def:m3:defi-credit-and-leverage:algo}{stablecoin algorithmique} UST a connu un \omterm{def:m3:blockchains-for-traders:stable}{décrochage} (depeg) en mai et est tombé près de zéro (\cref{ch:m3:defi-credit-and-leverage})~; Three Arrows Capital, un fonds crypto dont on rapporte qu'il gère plus de
3~milliards de dollars qui avait emprunté des \omterm{def:m3:blockchains-for-traders:key}{jetons} et des dollars auprès de plusieurs prêteurs, a reçu de leur part des notifications de défaut à mesure que les cours baissaient, et a été placé en liquidation
aux îles Vierges britanniques le 27 juin~; FTX a fait faillite en novembre
(\cref{ch:m3:centralised-exchanges}).

\section{Ce que voit chaque desk}

Les mêmes mécanismes ne se présentent pas de la même façon selon le desk. Pour la société de trading, les leçons sont concrètes. Le collatéral est la contrainte~: dans chaque épisode, les
perdants n'avaient pas tort sur la valeur, mais ils manquaient de cash lorsque les marges ont augmenté, si bien que le plan de liquidité de la société doit supposer que les marges augmentent sur toutes ses plateformes
à la fois. Les couvertures échouent lorsque les corrélations se rompent~: une couverture estimée en période calme (obligations contre actions, bitcoin contre actions, une plateforme contre une autre) peut évoluer dans le même sens
lors d'une ruée vers la liquidité (dash for cash). Des détenteurs communs rapprochent des marchés éloignés~: la société doit savoir qui d'autre détient ce qu'elle détient, et avec quel effet de levier.
Enfin, les tests de résistance (livre~6) doivent inclure le scénario conjoint, et non seulement la pire journée de chaque marché. Le simulateur de ce chapitre illustre ce point par des
chiffres.

\begin{omfigure}
\begin{tikzpicture}
\begin{axis}[width=11cm,height=5cm,
  xlabel={jour},ylabel={prix},
  tick label style={font=\small},label style={font=\small},
  xmin=0,xmax=39,ymin=70,ymax=110,grid=major,grid style={gray!25},
  legend columns=3,legend style={font=\footnotesize,at={(0.5,-0.3)},anchor=north,draw=none,fill=none,/tikz/every even column/.append style={column sep=8pt}},
  table/col sep=comma]
\addplot[omThm,very thick] table[x=day,y=a1]{figdata/markets-3/28-when-markets-break-together/paths.csv};\addlegendentry{actif 1 (choqué)}
\addplot[omDef,very thick,dashed] table[x=day,y=a2]{figdata/markets-3/28-when-markets-break-together/paths.csv};\addlegendentry{actif 2}
\addplot[omExo,very thick,dotted] table[x=day,y=a3]{figdata/markets-3/28-when-markets-break-together/paths.csv};\addlegendentry{actif 3}
\draw[gray,thick] (axis cs:10,70) -- (axis cs:10,110);
\end{axis}
\end{tikzpicture}

\omcaption{La spirale de marge du tutoriel~: 40 détenteurs avec un effet de levier de 8 fois sur trois actifs~; au jour 10, l'actif 1 chute de 10~\%. La hausse des marges et
la baisse des fonds propres forcent la vente des trois actifs, et les actifs 2 et 3, non touchés par le choc, baissent avec lui~; la corrélation des actifs 1 et 2 passe de 0.30 dans les dix jours précédant le choc à
0.99 dans les dix jours qui suivent, la baisse conjointe du jour du choc dominant. Données synthétiques. Données~: le tutoriel du chapitre.}\label{fig:m3:when-markets-break-together:paths}
\end{omfigure}

\begin{table}[ht]
\centering
\small
\begin{tabular}{@{}lrrr@{}}
\toprule
Cas & \omterm{def:m3:when-markets-break-together:fire}{ventes forcées}, Md\$ & multiplicateur & corrélation après \\
\midrule
Cas de base & 223.6 & 8.88 & 0.99 \\
Huit pas par jour & 225.6 & 8.73 & 0.99 \\
Sans \omterm{def:m3:when-markets-break-together:spirals}{spirale des marges} (marges fixes) & 0.0 & 0.00 & 0.11 \\
Sans \omterm{def:m3:when-markets-break-together:spirals}{spirale des pertes} (capital non réévalué) & 121.9 & 4.94 & 0.98 \\
Sans impact sur les prix & 124.4 & 4.94 & 0.11 \\
Un actif par détenteur & 183.8 & 4.67 & 0.04 \\
\bottomrule
\end{tabular}
\caption{La spirale et ses variantes d'ablation~: \omterm{def:m3:when-markets-break-together:fire}{ventes forcées} sur l'ensemble de la simulation, \omterm{def:m3:when-markets-break-together:fire}{ventes forcées} par dollar de déficit de marge créé par le choc (le
multiplicateur de \omterm{def:m3:when-markets-break-together:fire}{ventes forcées}), et corrélation des actifs 1 et 2 sur les dix jours suivant le choc. Sans impact sur les prix, le multiplicateur vaut le $1/m$ de la
\cref{prop:m3:when-markets-break-together:multiplier}~; avec impact, il est presque doublé~; un pas de temps plus fin le modifie peu~; avec des marges fixes, ce choc est
absorbé sans \omterm{def:m3:when-markets-break-together:fire}{ventes forcées}.}\label{tab:m3:when-markets-break-together:ablations}
\end{table}

\section{Tutoriel~: une spirale de marge sur trois actifs}\label{tut:m3:when-markets-break-together:spiral}

\begin{tutorial}
\textbf{Objectif.} Simuler une spirale de marge sur trois actifs détenus par des investisseurs à effet de levier, les marges augmentant avec la volatilité et les \omterm{def:m3:when-markets-break-together:fire}{ventes forcées} faisant bouger les prix,
mesurer la hausse des corrélations lors du stress, vérifier le résultat avec un pas de temps plus fin, et désactiver chaque mécanisme tour à tour. \textbf{État final~:}
\cref{fig:m3:when-markets-break-together:paths}, \cref{tab:m3:when-markets-break-together:ablations} et les chiffres du problème du week-end.
\begin{tutsteps}
\item \textbf{Désendettement.} Chaque détenteur en déficit de marge vend la même fraction de tout ce qu'il détient~; les ventes font bouger les prix.
\omcode{code/firm/marginspiral/firm_marginspiral.py}{100}{119}{Ventes forcées et leur impact sur les prix.}\label{lst:m3:when-markets-break-together:delever}
\item \textbf{Stabilité et ablations.} \texttt{ablation\_table()} exécute le cas de base, huit pas par jour, et chaque mécanisme désactivé tour à tour.
\item \textbf{Exécuter} \texttt{correlations()} et \texttt{fig\_break.py}.
\end{tutsteps}
\textbf{Pour aller plus loin.} Donner aux détenteurs des effets de levier différents et voir qui fait défaut en premier~; ajouter un coussin de liquidités et trouver la taille à partir de laquelle la spirale s'arrête~;
laisser les marges réagir avec un retard, comme le ferait le modèle d'une chambre de compensation fondé sur une fenêtre d'observation passée.
\end{tutorial}

\section{Construction~: le simulateur de spirale de marge}\label{bld:m3:when-markets-break-together:marginspiral}

\begin{build}
\bfield{Rôle} Les tests de résistance de la petite société doivent inclure le scénario conjoint dans lequel ses propres marges augmentent sur toutes les plateformes pendant que les actifs qu'elle détient baissent
ensemble~; le simulateur produit ces scénarios et mesure l'ampleur des \omterm{def:m3:when-markets-break-together:fire}{ventes forcées} que crée un choc.
\bfield{Interface} \texttt{Params} (volatilités, corrélation, marges de base, impact, détenteurs, effet de levier, jours, choc, décroissance EWMA, sous-pas, interrupteurs pour la
spirale de marge, la spirale de pertes, la détention croisée et les \omterm{def:m3:when-markets-break-together:fire}{ventes forcées})~; \texttt{simulate(params)} renvoyant les prix, les rendements, les \omterm{def:m3:when-markets-break-together:fire}{ventes forcées}, le déficit initial et
les défauts~; \texttt{forced\_sale\_multiplier}~; \texttt{correlation(result, a, b, days)}.
\bfield{Règles} Le désendettement itéré jusqu'à un point fixe à l'intérieur de chaque pas, de sorte que les résultats ne dépendent pas de la taille du pas~; impact logarithmique sur les prix, de sorte que les prix restent
positifs~; chaque mécanisme peut être désactivé.
\bfield{Tests d'acceptation} \texttt{code/firm/marginspiral/tests/}~: le multiplicateur vaut $1/m$ sans impact ni spirale de pertes~; stabilité avec un pas de temps plus fin~;
chaque ablation réduit les \omterm{def:m3:when-markets-break-together:fire}{ventes forcées} ou la corrélation~; aucune \omterm{def:m3:when-markets-break-together:fire}{vente forcée} lorsque les \omterm{def:m3:when-markets-break-together:fire}{ventes forcées} sont désactivées.
\bfield{Pour aller plus loin} Détenteurs hétérogènes et coussin de liquidités~; modèles de marge décalés dans le temps~; un second choc~; calibrage de l'impact sur la profondeur d'un marché.
\end{build}

\begin{omsources}
\item M.~K.~Brunnermeier et L.~H.~Pedersen, ``Market Liquidity and Funding Liquidity'', \emph{Review of Financial Studies} 22(6), 2009.
\item Financial Stability Board, \emph{Holistic Review of the March Market Turmoil}, 17 novembre 2020.
\item BCBS--CPMI--IOSCO, communiqué de presse sur les pratiques de marge, 29 septembre 2022~; communiqué de presse de LCH.Clearnet sur le défaut de Lehman, 8 octobre 2008 (voir les registres du
livre~1).
\item FRED, séries DGS10, DTWEXBGS et CBBTCUSD~; chapitres 8, 13, 15 et 23 de ce livre et le livre~2, chapitre~7, pour 2022.
\item BCE, \emph{Financial Stability Review}, novembre 2022, encadré spécial ``Financial stability risks from energy derivatives markets''; Fortum, communiqué de
presse sur la facilité KfW d'Uniper, 29 août 2022.
\item In re Three Arrows Capital Ltd., No.\ 22-10920 (Bankr.\ S.D.N.Y.), verified petition and declaration of the joint liquidators, 1 juillet 2022.
\end{omsources}

\section{Exercices}

\begin{exercise}[$\star$]\label{exo:m3:when-markets-break-together:1}
Un détenteur est en déficit de marge de 10 millions de dollars sur des positions margées à 12.5~\%. De combien doit-il vendre si les prix ne bougent pas~?
\end{exercise}

\begin{exercise}[$\star$]\label{exo:m3:when-markets-break-together:2}
Distinguer une \omterm{def:m3:when-markets-break-together:flight}{fuite vers la qualité} d'une ruée vers la liquidité (dash for cash), et dire ce que chacune fait aux prix des obligations d'État.
\end{exercise}

\begin{exercise}[$\star$]\label{exo:m3:when-markets-break-together:3}
Pourquoi un actif considéré comme une couverture peut-il chuter avec les actions dans une ruée vers la liquidité (dash for cash)~?
\end{exercise}

\begin{exercise}[$\star\star$]\label{exo:m3:when-markets-break-together:4}
D'après le \cref{tab:m3:when-markets-break-together:ablations}, quelle part de la \omterm{def:m3:when-markets-break-together:fire}{vente forcée} du cas de base est due, prises ensemble, à l'impact sur les prix et à la \omterm{def:m3:when-markets-break-together:spirals}{spirale des pertes}~?
\end{exercise}

\begin{exercise}[$\star\star$]\label{exo:m3:when-markets-break-together:5}
Le rendement à dix ans est passé de 0.54~\% à 1.18~\% entre le 9 et le 18 mars 2020. Avec une duration modifiée de 9, combien a perdu un bon du Trésor à dix ans~?
\end{exercise}

\begin{exercise}[$\star\star$]\label{exo:m3:when-markets-break-together:6}
Pourquoi une \omterm{def:m3:when-markets-break-together:flight}{rupture des corrélations} rend-elle invalide une valeur en risque estimée sur des données calmes~?
\end{exercise}

\begin{exercise}[$\star\star\star$]\label{exo:m3:when-markets-break-together:7}
\emph{Programmation.} Avec \texttt{run}, diviser le levier par deux. Que deviennent les \omterm{def:m3:when-markets-break-together:fire}{ventes forcées} et le multiplicateur~?
\end{exercise}

\begin{exercise}[$\star\star\star$]\label{exo:m3:when-markets-break-together:8}
\emph{Repérer l'erreur.} « Notre portefeuille est diversifié entre taux, matières premières et crypto-actifs, donc une crise dans l'un n'atteindra pas les autres. »
\end{exercise}

\section{Problème~: mars 2020 à travers les desks}

\begin{problem}[{Problème du week-end --- combien de vente une choc force-t-il ?}]\label{pb:m3:when-markets-break-together:1}
Utiliser le simulateur du tutoriel~: 40 détenteurs, chacun avec 1 milliard de dollars de fonds propres et un effet de levier de 8 réparti sur trois actifs, des marges qui augmentent avec la volatilité EWMA, et
un choc de 10~\% sur l'actif 1 le jour 10.

\medskip\noindent\textbf{Partie I --- Le choc.}
\begin{enumerate}
\item Quel déficit de marge le choc crée-t-il chez les détenteurs~?
\item Combien vendent-ils de force au cours de la simulation~?
\item Quel est le multiplicateur de \omterm{def:m3:when-markets-break-together:fire}{vente forcée}~?
\item Quel serait le multiplicateur sans impact sur les prix, et pourquoi~?
\item Que se passe-t-il pour les actifs 2 et 3, que le choc n'a pas touchés~?
\end{enumerate}

\medskip\noindent\textbf{Partie II --- Les mécanismes.}
\begin{enumerate}[resume]
\item Que devient le résultat avec huit pas par jour, et pourquoi cela compte-t-il~?
\item Que se passe-t-il avec des marges fixes~?
\item Et si chaque détenteur ne possède qu'un seul actif~?
\item Comment la corrélation entre les actifs 1 et 2 change-t-elle après le choc~?
\item Quel mécanisme produit l'essentiel de l'amplification~?
\end{enumerate}

\medskip\noindent\textbf{Partie III --- L'histoire.}
\begin{enumerate}[resume]
\item Que le desk de taux a-t-il vu en mars 2020~?
\item Que firent les chambres de compensation des marges ce mois-là~?
\item Pourquoi le traitement de Lehman par LCH en 2008 n'a-t-il pas empêché la crise~?
\item Quel épisode de 2022 ressemble le plus au simulateur, et pourquoi~?
\item Qu'ont fait les banques centrales dans chaque épisode que le simulateur laisse de côté~?
\end{enumerate}

\medskip\noindent\textbf{Partie IV --- Jugement.}
\begin{enumerate}[resume]
\item Les chambres de compensation doivent-elles relever les marges en période de crise~?
\item Comment une société de trading doit-elle dimensionner son coussin de liquidité pour un stress conjoint~?
\item Que surveilleriez-vous pour voir venir une spirale~?
\item Énoncer le \emph{résultat nommé}~: le multiplicateur de \omterm{def:m3:when-markets-break-together:fire}{vente forcée} de la spirale pour le levier et l'impact indiqués, et sans impact.
\item En une phrase~: pourquoi les marchés se brisent-ils ensemble~?
\end{enumerate}
\end{problem}

\section{Questions d'entretien}

\begin{interviewq}[$\star$ \iqroles{risque}]\label{iq:m3:when-markets-break-together:1}
Qu'est-ce qu'une \omterm{def:m3:when-markets-break-together:spirals}{spirale des pertes} et une \omterm{def:m3:when-markets-break-together:spirals}{spirale des marges}~?
\end{interviewq}

\begin{interviewq}[$\star$ \iqroles{trader}]\label{iq:m3:when-markets-break-together:2}
Pourquoi les bons du Trésor ont-ils chuté au milieu de mars 2020~?
\end{interviewq}

\begin{interviewq}[$\star\star$ \iqroles{chercheur}]\label{iq:m3:when-markets-break-together:3}
Comment estimeriez-vous les corrélations pour un scénario de stress~?
\end{interviewq}

\begin{interviewq}[$\star\star$ \iqroles{risque}]\label{iq:m3:when-markets-break-together:4}
Concevez un test de stress de liquidité pour une société qui négocie des contrats à terme, des perpétuels crypto et des contrats de change à terme.
\end{interviewq}

\begin{interviewq}[$\star\star$ \iqroles{trader, risque}]\label{iq:m3:when-markets-break-together:5}
Votre client, un hedge fund, est long de bons du Trésor et vendeur de contrats à terme avec un effet de levier de 50. Que lui arrive-t-il dans une ruée vers la liquidité (dash for cash)~?
\end{interviewq}

\begin{interviewq}[$\star\star\star$ \iqroles{chercheur, développeur}]\label{iq:m3:when-markets-break-together:6}
Construisez un simulateur de \omterm{def:m3:when-markets-break-together:spirals}{spirale des marges} et expliquez comment vous vous assureriez que ses résultats ne sont pas un artefact de son pas de temps.
\end{interviewq}
